\section{Introduction}
\label{sec:intro}

Synthesizing high-fidelity novel view images of a 3D scene is imperative for various industrial applications, ranging from gaming and filming to AR/VR. In recent years, Neural Radiance Fields(NeRF) ~\cite{NeRF} has demonstrated its remarkable ability on this task with photorealistic renderings. While lots of subsequent works focus on improving rendering quality~\cite{Mip-NeRF, Mip-NeRF360} or training and rendering speed~\cite{Instant-NGP, TensoRF, EfficientNeRF, Plenoxels} for static scenes, another line of work~\cite{D-NeRF, KPlanes, TiNeuVox} proposes to extend the setting to a dynamic scene. Though various attempts have been made to improve efficiency and dynamic rendering quality, the introduction of an additional time-variant MLP to model complex motions in dynamic scenes will inevitably cause a surge in computational cost during both the training and inference process.

More recently, 3D Gaussian Splatting(3D-GS)~\cite{3D-GS} manages to achieve real-time rendering with high visual quality by introducing a novel point-like representation, referred to as 3D Gaussians. However, it mainly deals with static scenes. Though methods such as leveraging a deformation network on each 3D Gaussian can extend 3D-GS to dynamic scenes, the rendering speed will be greatly affected, especially when a large number of 3D Gaussians is necessary to represent the scene.

Drawing inspiration from the well-established \textit{As-Rigid-As-Possible} regularization in 3D reconstruction and the superpoint/superpixel concept in point clouds/images over-segmentation, we propose a novel approach named Superpoint Gaussian Splatting (SP-GS) in this work for reconstructing and rendering dynamic scenes. 
The key insight lies in that each 3D Gaussian should not be a completely independent entity. Some neighboring 3D Gaussians probably possess similar translation and rotation transformations at all timesteps due to the properties of a rigid motion. We can cluster these similar 3D Gaussians together to form a superpoint so that it is no longer necessary to compute a deformation for every single 3D Gaussian, leading to a much faster rendering speed.

To be specific, after acquiring the initial 3D Gaussians of the canonical space through a warm-up training process, a learnable association matrix will be applied to the initial 3D Gaussians and group them into several superpoints. Subsequently, our framework will leverage a tiny MLP network for predicting the deformations of superpoints, which will later be utilized to compute the deformation of every single 3D Gaussian in the superpoint to enable novel view rendering for dynamic scenes. Apart from the rendering loss at each timestep, to take full advantage of the 
 \textit{As-Rigid-As-Possible} feature within one superpoint, we additionally utilize a property reconstruction loss on the properties of Gaussians, including positions, translations, and rotations. 

Thanks to the computational expense saved by using superpoints, our approach manages to achieve a comparable rendering speed with 3D-GS. Furthermore, the mixed representation of 3D Gaussians and superpoints possess a strong extensibility like adding a non-rigid motion prediction module for better dynamic scene reconstruction. Last but not least, SP-GS can facilitate various downstream applications like editing a reconstructed scene as superpoints should cluster similar 3D Gaussians together, providing more meaningful groups than pure 3D Gaussians. 
Our contributions can be summarized as follows:
\begin{itemize}
  \item We introduce Superpoint Gaussian Splatting (SP-GS), a novel approach for high-fidelity and real-time rendering in dynamic scenes that aggregates 3D Gaussians with similar deformations into superpoints. 
  \item Our method possesses a strong extensibility like adding a non-rigid prediction module or distillation from a larger model and can facilitate various downstream applications like scene editing.
  \item SP-GS achieves real-time rendering on dynamic scenes, up to 227 FPS at a resolution of $800\times 800$ for synthetic datasets and 117 FPS at a resolution of $536 \times 960$ in real datasets with superior or comparable performance than previous SOTA methods.
\end{itemize}
% See our project page for video results and source code at \url{https://dnvtmf.github.io/SP_GS.github.io}.
